%%% Hlavní soubor. Zde se definují základní parametry a odkazuje se na ostatní části. %%%

%% Verze pro jednostranný tisk:
% Okraje: levý 40mm, pravý 25mm, horní a dolní 25mm
% (ale pozor, LaTeX si sám přidává 1in)
\documentclass[12pt,a4paper]{report}
\setlength\textwidth{145mm}
\setlength\textheight{247mm}
\setlength\oddsidemargin{15mm}
\setlength\evensidemargin{15mm}
\setlength\topmargin{0mm}
\setlength\headsep{0mm}
\setlength\headheight{0mm}
% \openright zařídí, aby následující text začínal na pravé straně knihy
\let\openright=\clearpage

%% Pokud tiskneme oboustranně:
% \documentclass[12pt,a4paper,twoside,openright]{report}
% \setlength\textwidth{145mm}
% \setlength\textheight{247mm}
% \setlength\oddsidemargin{15mm}
% \setlength\evensidemargin{0mm}
% \setlength\topmargin{0mm}
% \setlength\headsep{0mm}
% \setlength\headheight{0mm}
% \let\openright=\cleardoublepage

%% Použité kódování znaků: obvykle latin2, cp1250 nebo utf8:
\usepackage[utf8]{inputenc}

%% Ostatní balíčky
\usepackage{graphicx}
\usepackage{amsthm}
\usepackage{amssymb}
\usepackage{amsmath}
\usepackage{framed}
\usepackage[titles]{tocloft}

%% Balíček hyperref, kterým jdou vyrábět klikací odkazy v PDF,
%% ale hlavně ho používáme k uložení metadat do PDF (včetně obsahu).
%% POZOR, nezapomeňte vyplnit jméno práce a autora.
\usepackage[ps2pdf,unicode]{hyperref}   % Musí být za všemi ostatními balíčky
\hypersetup{pdftitle=Separation axioms}
\hypersetup{pdfauthor=Karel Ha}

%%% Drobné úpravy stylu

% Tato makra přesvědčují mírně ošklivým trikem LaTeX, aby hlavičky kapitol
% sázel příčetněji a nevynechával nad nimi spoustu místa. Směle ignorujte.
\makeatletter
\def\@makechapterhead#1{
  {\parindent \z@ \raggedright \normalfont
   \Huge\bfseries \thechapter. #1
   \par\nobreak
   \vskip 20\p@
}}
\def\@makeschapterhead#1{
  {\parindent \z@ \raggedright \normalfont
   \Huge\bfseries #1
   \par\nobreak
   \vskip 20\p@
}}
\makeatother

% Toto makro definuje kapitolu, která není očíslovaná, ale je uvedena v obsahu.
\def\chapwithtoc#1{
\chapter*{#1}
\addcontentsline{toc}{chapter}{#1}
}

%%%%%%%%%%%%%%%%%%%%%%%%%%%%%%%%%%%%%%%%%%%%%%
% přidal Karel Ha
\newcommand{\HRule}{\rule{\linewidth}{0.5mm}}
\theoremstyle{plain}
\newtheorem{thm}{Theorem}[section]
\newtheorem{lem}[thm]{Lemma}
\newtheorem{prop}[thm]{Proposition}
\newtheorem{cor}[thm]{Corollary}

\theoremstyle{definition}
\newtheorem*{df}{Definition}	% Definice nečíslujeme, proto "*"
\newtheorem{exmpl}{Example}
\newtheorem{fact}{Fact}
\newtheorem*{nota}{Notation}

\theoremstyle{remark}
\newtheorem{rem}[thm]{Remark}
\newtheorem{note}[thm]{Note}

\numberwithin{equation}{section}
\renewcommand{\theequation}{\arabic{chapter}.\alph{equation}}

\renewcommand{\qedsymbol}{$\boxtimes$}
%\renewcommand{\qedsymbol}{$\blacksquare$}

\def\labelitemi{$\spadesuit$}
\renewcommand\thechapter{\Roman{chapter}} 
%\renewcommand\thesection{\arabic{section}}
%\renewcommand\thesubsection{(\alph{subsection})}

\renewcommand{\labelenumi}{{\rm(\roman{enumi})}}

% mezera v Table of Contents před číslem kapitoly
\renewcommand*\cftchapnumwidth{2.4em}
\renewcommand*\cftsecnumwidth{3em}
% samotna cisla "section" misto "chapter.section"
%\renewcommand\cftchapaftersnum{}
%\renewcommand\cftsecaftersnum{.}

\newcommand{\N}{\mathbb{N}}
\newcommand{\Q}{\mathbb{Q}}
\newcommand{\R}{\mathbb{R}}
\newcommand{\I}{\mathbb{I}}
\newcommand{\Acal}{\mathcal{A}}
\newcommand{\Bcal}{\mathcal{B}}
\newcommand{\Ccal}{\mathcal{C}}
\newcommand{\D}{\mathfrak{D}}
\newcommand{\none}{\varnothing}
\newcommand{\st}{\mid}
\newcommand{\uparr}{\left\uparrow\right.}
\newcommand{\darr}{\left\downarrow\right.}
\newcommand{\dL}{\left\uparrow d_L \right.}
\newcommand{\n}{\mathfrak{n}}
\newcommand{\rb}{\prec}
\newcommand{\cb}{\prec\mkern-8mu\prec}
\renewcommand{\phi}{\varphi}

\hyphenation{approach}
\hyphenation{ca-te-go-ry}
\hyphenation{con-ti-nuous}
\hyphenation{de-fi-ning}
\hyphenation{equi-va-len-ce}
\hyphenation{equi-va-len-ces}
\hyphenation{pro-duct}
\hyphenation{pro-ducts}
\hyphenation{se-cond-ly}
\hyphenation{se-pa-ra-tion}
\hyphenation{strong-ly}
\hyphenation{sub-lo-ca-le}
\hyphenation{weak-ly}
%%%%%%%%%%%%%%%%%%%%%%%%%%%%%%%%%%%%%%%%%%%%%%

\begin{document}

% Trochu volnější nastavení dělení slov, než je default.
\lefthyphenmin=2
\righthyphenmin=2

%%% Titulní strana práce

\pagestyle{empty}
\begin{center}

\large

Charles University in Prague

\medskip

Faculty of Mathematics and Physics

\vfill

{\bf\Large BACHELOR THESIS}

\vfill

\centerline{\mbox{\includegraphics[width=60mm]{../img/logo.eps}}}

\vfill
\vspace{5mm}

{\LARGE Karel Ha}

\vspace{15mm}

% Název práce přesně podle zadání
\HRule \\[0.6cm]
{\LARGE\bfseries Separation axioms} \\[0.2cm]
\HRule \\[1.5cm]

\vfill

% Název katedry nebo ústavu, kde byla práce oficiálně zadána
% (dle Organizační struktury MFF UK)
Department of Applied Mathematics

\vfill

\begin{tabular}{rl}

Supervisor of the bachelor thesis: & prof. RNDr. Aleš Pultr, DrSc. \\
\noalign{\vspace{2mm}}
Study programme: & Computer Science \\
\noalign{\vspace{2mm}}
Specialization: & General Computer Science \\
\end{tabular}

\vfill

% Zde doplňte rok
Prague 2013

\end{center}

\newpage

%%% Následuje vevázaný list -- kopie podepsaného "Zadání bakalářské práce".
%%% Toto zadání NENÍ součástí elektronické verze práce, nescanovat.

%%% Na tomto místě mohou být napsána případná poděkování (vedoucímu práce,
%%% konzultantovi, tomu, kdo zapůjčil software, literaturu apod.)

\openright

\noindent
Dedication.

\newpage

%%% Strana s čestným prohlášením k bakalářské práci

\vglue 0pt plus 1fill

\noindent
I declare that I carried out this bachelor thesis independently, and only with the cited
sources, literature and other professional sources.

\medskip\noindent
I understand that my work relates to the rights and obligations under the Act No.
121/2000 Coll., the Copyright Act, as amended, in particular the fact that the Charles
University in Prague has the right to conclude a license agreement on the use of this
work as a school work pursuant to Section 60 paragraph 1 of the Copyright Act.

\vspace{10mm}

\hbox{\hbox to 0.5\hsize{%
In ........ date ............
\hss}\hbox to 0.5\hsize{%
signature of the author
\hss}}

\vspace{20mm}
\newpage

%%% Povinná informační strana bakalářské práce

\vbox to 0.5\vsize{
\setlength\parindent{0mm}
\setlength\parskip{5mm}

Název práce:
Oddělovací axiomy
% přesně dle zadání

Autor:
Karel Ha

Katedra:  % Případně Ústav:
Katedra aplikované matematiky
% dle Organizační struktury MFF UK

Vedoucí bakalářské práce:
prof. RNDr. Aleš Pultr, DrSc., Katedra aplikované matematiky
% dle Organizační struktury MFF UK, případně plný název pracoviště mimo MFF UK

Abstrakt:
% abstrakt v rozsahu 80-200 slov; nejedná se však o opis zadání bakalářské práce

Klíčová slova:
% 3 až 5 klíčových slov
bezbodová topologie, oddělovací axiomy, locale, frame

\vss}\nobreak\vbox to 0.49\vsize{
\setlength\parindent{0mm}
\setlength\parskip{5mm}

Title:
% přesný překlad názvu práce v angličtině
Separation axioms

Author:
Karel Ha

Department:
Department of Applied Mathematics
% dle Organizační struktury MFF UK v angličtině

Supervisor:
prof. RNDr. Aleš Pultr, DrSc., Department of Applied Mathematics
% dle Organizační struktury MFF UK, případně plný název pracoviště
% mimo MFF UK v angličtině

Abstract:
% abstrakt v rozsahu 80-200 slov v angličtině; nejedná se však o překlad
% zadání bakalářské práce

Keywords:
% 3 až 5 klíčových slov v angličtině
pointless topology, point-free topology, separation axioms, locale, frame

\vss}

\newpage

%%% Strana s automaticky generovaným obsahem bakalářské práce. U matematických
%%% prací je přípustné, aby seznam tabulek a zkratek, existují-li, byl umístěn
%%% na začátku práce, místo na jejím konci.

\openright
\pagestyle{plain}
\setcounter{page}{1}
\tableofcontents

%%% Jednotlivé kapitoly práce jsou pro přehlednost uloženy v samostatných souborech
\chapter*{Introduction}
\addcontentsline{toc}{chapter}{Introduction}

In classic topology, spaces are often specified by \emph{separation axioms\/}
of~various strength.
The stronger axiom we have, the more ``geometrical'' a~space appears to be.

One of~the best-known series is the family of~$T$-axioms.
They concern the separation of~points, points from closed sets, and closed sets
from other closed sets.
In~this context the role of~points is fundamental:
it may seem immensely difficult to~imitate $T$-axioms in~point-free topology.
Nevertheless, most of them do have pointless counterparts and this thesis deals
with them.

\section*{Structure of thesis}

The first chapter provides basics of~order and category theories later followed
by~essentials of~point-set and point-free topologies.

The subsequent chapter presents \emph{subfitness\/}---a~weakened form of~the
$T_1$-axiom, which itself is not interesting in the point-free context, whereas
the subfitness is crucial for relations between other axioms.

Axioms of~\emph{Hausdorff type\/} have a~distinctive role:
unlike the others, we are not able to directly translate them from the
$T_2$-axiom.
For various reasons point-free versions cannot correspond to the point-set
Hausdorff property.

Chapters V and VI regard \emph{regularity\/} and \emph{complete regularity\/}.
We show how to exclude points from their definitions via ``rather-below''
relation $\rb$, and thus, how to create their point-free analogies.

The final chapter regards \emph{normality\/}, which we can very naturally
translate into the point-free language.

\chapter*{Preliminaries}
\addcontentsline{toc}{chapter}{Preliminaries}

\section*{Notations}

\begin{itemize}
\item $2^X$ is the set of~all subsets, explicitly, $2^X = \{ S \st S\subseteq X
\}$.

\item $A \Subset B$ means {\sl set $A$ is a finite subset of set $B$\/}.

\item $\Omega(X)$ denotes the lattice (or more specifically, the frame) of open
sets in $X$ together with $\subseteq$ as the ordering. 
By abuse of notation, we denote the related underlying set in the very same
way.

\item $\overline{S}$ is the~closure of $S$ in topology $(X, \Omega(X))$.
It is defined by
\[
  \overline{S} = \bigcap_{ \text{closed } A\colon \, A\supseteq S} A
  \, ,
\]
or equivalently,
\[
  \overline{S} = \{ x \in X \st U\in\Omega(X) \text{ and } U\owns x \Rightarrow
  U \cap S \ne \none \}.
\]
\item $f^*$ resp. $f_*$ is the~left resp. right adjoint of $f$ (viz ``Order
theory'').
\item $\left\downarrow U \right. = \bigcup_{u\in U} \{x \st x \le u\}$ and
$\left\uparrow U \right. = \bigcup_{u\in U} \{x \st x \ge u\}$.

\end{itemize}

\section*{Definitions}

\subsection*{Order theory}

\begin{itemize}
\item A \emph{pseudocomplement} of an element $x$ is $x^*$ such that $y \le x^*
\equiv x \wedge y = 0$.
\end{itemize}

\begin{fact}
  $x \le x^{**}$.
  (As the $x$, among others, meets with $x^*$ in $x^* \wedge x = 0$.)
\end{fact}

\begin{itemize}
\item A relation $R$ \emph{interpolates} if $R \subseteq R \circ R$.

\item Monotone maps $l\colon A \to B$ and $r\colon B \to A$ are \emph{Galois
  adjoints\/} ($l$ is a~\emph{left adjoint\/} of $r$; $r$ is a~\emph{right
  adjoint\/} of $l$) if for every $a\in A, \, b\in B$
  \[
    l(a) \le b \; \equiv a \le r(b).
  \]

\item A~set $U$ is a~\emph{down-set\/} resp. an~\emph{up-set\/} if
$\left\downarrow U \right. = U$ resp. $\left\uparrow U \right. = U$.
\end{itemize}

\subsection*{Category theory}

\begin{itemize}
\item A \emph{category} $\Ccal$ is a structure consisting of
  \begin{itemize}
  \item class~$\Ccal^O$ of \emph{objects\/}
  \item class~$||\Ccal||$ of \emph{morphisms\/} between objects 
  (a~concrete morphism $f$ between objects $A$ and $B$ is denoted by $f\colon A
   \to B$)
  \end{itemize}
further satisfying axioms of
  \begin{itemize}
  \item \emph{(identity)\/}
  For every $X\in \Ccal^O$ there is a~morphism $1_X$ such that for any $f\colon A \to B$
  we have $1_A \cdot f = f = f \cdot 1_B$.
  \item \emph{(asociativity)\/}
  Whenever $f\colon A \to B$, $g\colon B \to C$ and $h\colon C \to D$ then $h
  \cdot (g \cdot f) = (h \cdot g) \cdot f$.
  \end{itemize}

\item The \emph{opposite category}~$\mathcal{C}^{op}$ of a category
$\mathcal{C}$ comprises the same objects $\Ccal^O$ and ``reverse'' morphisms
$f^{op}\colon A \leftarrow B$ (for each morphism $f\colon A \to B$
in~$||\Ccal||$).

\item An \emph{epimorphism} is a right-cancellative morphism $e$, that is,
\[
  f_1 \cdot e = f_2 \cdot e \; \Longrightarrow \; f_1 = f_2.
\]
\end{itemize}

\begin{exmpl}
  In the category {\bf Loc} (see ``Point-free topology'') every localic map
  onto is epic.
\end{exmpl}

\begin{itemize}
\item A \emph{product} of a system $\left(A_i\right)_{i\in J}$ is a~system
of~\emph{projections\/}
\[
  \left(p_i\colon \prod_{j\in J} A_j \to A_i \right)_{i\in J}
\]
with the property that for arbitrary system $\left(f_i\colon X \to
A_i\right)_{i\in J}$ there is a~unique solution~$f\colon X \to \prod_{j\in
J} A_j$ to equations
\[
  p_i \cdot f = f_i \text{ for } i\in J.
\]

\item Similarly, a \emph{coproduct} is a product in the opposite category.
Namely, it is a~system of~\emph{injections\/}
\[
  \left(\iota_i\colon A_i \to \coprod_{j\in J} A_j \right)_{i\in J}
\]
such that every $\left(g_i\colon A_i \to X \right)_{i\in J}$ has a~single
solution~$g\colon \coprod_{j\in J} A_j \to X$ to~equations
\[
  g \cdot \iota_i = g_i \text{ for } i\in J.
\]

\item A \emph{diagonal} is (the only) morphism $\Delta\colon A \to \prod_{i\in
J} A$ solving the system $p_i\cdot \Delta = 1_A$ for $i \in J$.
Dually, a \emph{codiagonal} is the solution $\nabla\colon \coprod_{i\in J} A
\to A$ to the system $\nabla\cdot \iota_i = 1_A, \, i \in J$.
\end{itemize}


\subsection*{Point-set topology}

\begin{itemize}
\item A \emph{topological space} is a~structure of
  \begin{itemize}
  \item the set of \emph{points\/} $X$
  \item \emph{open sets\/} $\tau\subseteq 2^X$ holding axioms
    \begin{description}
    \item[(o1)] $\none, X \in \tau$
    \item[(o2)] $\left\{ U_i \st i\in J \right\} \subseteq \tau \; \Rightarrow
    \; \bigcup_{i\in J} U_i \in \tau$
    \item[(o3)] $U_1, U_2 \in \tau \; \Rightarrow \; U_1 \cap U_2 \in \tau$
    \end{description}
  \end{itemize}

\item The \emph{basis} of a topological space $(X, \tau)$ is the
$\mathcal{B}\subseteq \tau$ such that for every open $U$ we can write $U =
\bigcup \{V\in \mathcal{B} \st V\subseteq U \}$.

\item Likewise, $\mathcal{S}\subseteq \tau$ is a~\emph{subbasis} of a
topological space $(X, \tau)$ if $\left\{ \bigcap \mathcal{F} \st \mathcal{F}
\Subset \mathcal{S} \right\}$ is basis of~$\tau$.
\end{itemize}

\begin{exmpl}
  The topology $\I$ consists of underlying set $\langle 0, 1 \rangle$ and open
  sets generated by subbasis $\{ \langle 0, a \langle, \; \rangle a, 1 \rangle
  \st a\in \langle 0, 1 \rangle \}$.
\end{exmpl}

\begin{fact}
  The topological space $\left( \prod_{i\in J} X_i, \mathcal{S} \right)$, where
  $\mathcal{S} = \left\{ p_{i}^{-1}[U] \st i\in J, \; U\in \tau_i \right\}$
  denotes its subbasis, is the product of topological spaces $\left( X_i
  \right)_{i\in J}$;
  more precisely, the projections $\left( p_i\colon \left(\prod_{i\in J} X_i,
  \mathcal{S}\right) \to \left(X_i, \tau_i\right) \right)_{i\in J}$
  constitute the product in~the~category~{\bf Top}.
\end{fact}

\begin{itemize}
\item A function $f\colon X \to Y$ between topological space $X$ and $Y$ is
\emph{continuous} if 
\[
  U\in \Omega(Y) \; \Longrightarrow \; f^{-1}[U]\in \Omega(X).
\]
Since preimage preserves unions and finite intersections, this amounts to
requiring that $f^{-1}[U]$ is open for {\bf subbasic} $U$.
\end{itemize}

\subsection*{Point-free topology}

\begin{itemize}
\item A \emph{frame} is a~complete lattice $L$ satisfying ``infinite'' version
of distributivity
\[
  \bigvee A \wedge b = \bigvee \{ a \wedge b \st a\in A \}
\]
for every $A\subseteq L$ and $b\in L$.

\item A \emph{frame homomorphism} (between frames $M$ and $L$) is a~map
$h\colon M \to L$ preserving all joins (in particular, $h(0) = 0$) and~all
finite meets (especially, $h(1) = 1$).

\item The category of~frames and frame homomorphisms will be denoted by {\bf
Frm}.

\item A \emph{locale} is a~frame (it has different name because of~different
subsequent morphisms).

\item A \emph{localic map} (between locales $L$ and $M$) is a mapping $f\colon
L \to M$ having left (Galois) adjoints $f^*\colon M \to L$, which is also a
frame homomorphism.
Hence, localic maps are infima-preserving $f$ and their left adjoints $f^*$
preserves finite meets.

\item We will write {\bf Loc} for the category of~locales and localic maps.
Thus, $\mathbf{Loc} = \mathbf{Frm}^{op}$.

\item A \emph{sublocale} of locale $L$ is a $S\subseteq L$ satisfying
  \begin{description}
  \item[(S1)] $M\subseteq S \Longrightarrow \bigwedge M\in S$
  \item[(S2)] $x\in L \text{ and } s\in S \Longrightarrow x \rightarrow s \in S$ 
  \end{description}
\end{itemize}

\begin{fact}
  The set-theoretic image $f[L]$ under a localic map $f\colon L\to M$
  is a~sublocale of~$M$.
\end{fact}

\subsubsection*{On coproducts in Frm}

In the category of~frames we have a~coproduct
\[
  \iota_i: L_i \to L_1 \oplus L_2, \, i = 1, 2
\]
(there is a general coproduct but we will need a coproduct of two objects
 only).
It can be constructed as follows.

\begin{itemize}
\item A down-set $U\subseteq L_1 \times L_2$ is \emph{saturated\/} if for all
systems $(x_i)_{i\in J}$ in $L_1$ resp. $L_2$ and all $y$ in $L_2$ resp. $L_1$
we have \label{df:satur}
  \begin{align*}
    \{ (x_i, y) \st i\in J \} \subseteq U \, &\Rightarrow \, (\bigvee x_i, y)\in U,
    \text{ and} \\
    \{ (y, x_i) \st i\in J \} \subseteq U \, &\Rightarrow \, (y, \bigvee x_i)\in U.
  \end{align*}

\item $L_1 \oplus L_2$ will be the set of~all saturated downsets in $L_1
\times L_2$.

\item The definition above concerns also the void $J$, and hence, for each
saturated down-set~$U$ we get
\[
  U\supseteq \n := \{ (x, y) \st x = 0 \text{ or } y = 0 \}.
\]
Since $\n$ itself is obviously saturated, it is \emph{the least saturated
down-set\/}.

\item For $(a, b)\in L_1 \times L_2$ the down-set
\[
  \darr (a, b) \cup \n
\]
is saturated.
This element of $L_1 \oplus L_2$ is denoted by
\[
  a \oplus b
\]
and we have the coproduct injections $\iota_i: L_i \to L_1 \oplus L_2$ defined
by
\[
  \iota_1(a) := a \oplus 1, \; \iota_2(b) := 1 \oplus b.
\]

\begin{cor*}
  Due to the coordinatewise meet, we have
  \[
    a \oplus b = \iota_1(a) \wedge \iota_2(b).
  \]
\end{cor*}

\item Note that for all saturated $U$ we have
\[
  U
  = \bigcup \{ (a, b) \st (a, b)\in U \}
  = \bigcup \{ a \oplus b \st (a, b)\in U \}
\]
(the second equality because $U$ is a~down-set).
As the set-theoretic union of down-sets is also a~down-set, it coincides with
their join.
Therefore, we may write 
\[
  U = \bigvee \{ a \oplus b \st (a, b)\in U \},
\]
and thus, \emph{the elements $a \oplus b$ generate $L_1 \oplus L_2$\/}.
\end{itemize}

For full details consult Chapter~IV of Picado and Pultr \cite{picado-pultr12}.

We have the following facts \ldots

\clearpage

\begin{itemize}
\item A \emph{closed sublocale} $S$ is a \ldots
Besides, the \emph{closure} of $S$ is the least closed sublocale containing $S$.
Since \ldots On the other hand, \ldots
\end{itemize}

\begin{exmpl} \label{(co)diag-in-Loc}
  The diagonal localic map $\Delta\colon L \to L \oplus L$ in {\bf Loc} is
  given by formula
  \[
    \Delta(a) = \left\{ (x, y) \st x \wedge y \leq a \right\},
  \]
  and likewise, its left adjoint (the codiagonal $\nabla\colon L \oplus L \to
  L$) is described by
  \[
    \nabla(U) = \bigvee \left\{ x \st (x, x)\in U \right\},
  \]
  as it is more explained in Corollary~IV.5.5.2 and Section~IV.5.2 of~Picado
  and Pultr \cite{picado-pultr12}.
\end{exmpl}

\section*{Convention}

In the whole text we will neglect the $T_0$ axiom, that is, we will assume all
topological spaces to be $T_0$-spaces.
Thus, every topological space $(X, \tau)$ satisfies the condition
\begin{center} \it
  for any $x \ne y$ from $X$ there is an open set $U \in \tau$ such that $x \in U
  \not\owns y$ or $y \in U \not\owns x$.
\end{center}

\chapter{Weakly and strongly Hausdorff}

In this chapter we will contemplate point-free analogies of the $T_2$
axiom---a~separation perhaps intricate to imitate in frames.
Rather, two resembling conditions will be presented.

\section{The $T_2$ axiom}

\begin{framed}
  \begin{df}[$T_2$]
    A topological space $(X, \tau)$ is called \emph{Hausdorff\/} (or, as well,
    the~\emph{$T_2$-space\/}) whenever
    \begin{center} \it
      for any $x \ne y$ from~$X$ there are open $U\owns x$ and $V\owns y$
      disjoint from~each other.
    \end{center}
  \end{df}
\end{framed}

\begin{rem} \label{T2->T1}
  Trivially from the definitions every $T_2$-space is, additionally, $T_1$.
  Thus, the sets of the form $X\setminus \left\{x\right\}$ are among the
  elements of $\Omega(X)$ (viz \ref{T1Char}).
\end{rem}

\section{The Dowker-Strauss approach}

In 1972 Dowker and Strauss \cite{ds72} introduced the condition:

\begin{framed}
  \begin{df}[$S_2'$]
    A locale $L$ satisfies the $S_2'$ axiom if for any $a, b \ne 1$ we have
    \[
      a \vee b = 1 \qquad \Rightarrow \qquad \exists u\not\leq a, v\not\leq
      b: \quad u \wedge v = 0.
    \]
  \end{df}
\end{framed}

As we are about to discover, it extends the Hausdorff property.

\begin{prop} \label{Haus->S2'}
  For every Hausdorff space $X$ the corresponding locale $\Omega(X)$ is always
  $S_2'$.
\end{prop}

\begin{proof}
  Open sets $A, B \ne X$ with $A \cup B = X$ contain $x\in X\setminus A =
  B\setminus A$ and $y\in X\setminus B = A\setminus B$.
  Particularly, we observe $x \ne y$ and, on top of that, receive the disjoint
  $U$ and $V$ from the $T_2$-definition.
  Evidently, $U\not\subseteq A$ (because of $x$) and $V\not\subseteq B$
  (because of $y$).
\end{proof}

\begin{rem}
  The reverse implication requires, in addition, the subfitness.
  For the~proof seek Proposition 7 of~Dowker and Strauss
  \cite{ds72}.
\end{rem}

For our purposes, nevertheless, $S_2'$ will be replaced by a~somewhat weaker
description:

\begin{framed}
  \begin{df}[DS-Haus]
    A locale $L$ is considered to be \emph{weakly Hausdorff\/} (also
    \emph{DS-Hausdorff} as in \emph{Dowker-Straus-Hausdorff}) if
    \[
      a \vee b \not\in \left\{a, b\right\} \qquad \Rightarrow \qquad \exists
      u\not\leq a, v\not\leq b: \quad u \wedge v = 0.
    \]
  \end{df}
\end{framed}

\begin{rem} \label{DSHaus->Haus}
  The analogous observation as in~\ref{Haus->S2'} holds:
  the~premise
  \[
    A \cup B \not\in \left\{ A, B \right\}
  \]
  is translated to
  \[
    A \not\subseteq B \textrm{ and } B \not\subseteq A;
  \]
  the rest of the proof can easily be finished in the same way.
\end{rem}

\section{Isbell's adaptation}

For the sake of future clarity, we now reveal another characterization of
$T_2$:

\begin{prop}
  A topological space $X$ is Hausdorff space if and only if the~diagonal
  $\left\{(x, x) \st x\in X \right\}$ is closed.
\end{prop}

\begin{proof}
  Let $\Delta$ stand for the mentioned diagonal.

  $\Rightarrow:$ Any element $(x, y)\not\in \Delta$, that means $x \ne y$, is
  separable from $\Delta$:
  namely, by~the~open $p_1^{-1}[U] \cap p_2^{-1}[V]$ with non-intersecting $U$
  and $V$ from the definition of~$(T_2)$.
  Thence, $(x, y)\not\in \overline{\Delta}$ concludes to $\Delta =
  \overline{\Delta}$, which is closed.

  $\Leftarrow:$ Choose $x \ne y$.
  In such a case, we get a $W\owns (x, y)$ from $\Omega(X\times X)$ fulfilling
  $W \cap \Delta = \none$.
  Without loss of generality, let $W$ be from basis.
  That is, $W$ is a~finite intersection
  \[
    \bigcap_{i=1}^m p_1^{-1}[U_i] \cap \bigcap_{j=1}^n p_2^{-1}[V_i] =
    p_1^{-1}[\bigcap_{i=1}^m U_i] \cap p_2^{-1}[\bigcap_{j=1}^n V_i],
  \]
  henceforth denoted only by $p_1^{-1}[U] \cap p_2^{-1}[V]$.

  Moreover, it is not possible for $W$ to be of the plain form $p_i^{-1}[U]$
  for $i = 1, 2$;
  since otherwise $p_i^{-1}[U] \cap \Delta \ne \none$, a contradiction.
  Following $W \cap \Delta = \none$, the open sets $U\owns x$ and $V\owns y$
  are disjoint from each other.
\end{proof}

Most likely, the just presented equivalence was the motivation for Isbell's
point-free approach to the $T_2$ axiom:
In~\cite{isbell72}, Hausdorff locales are exactly the ones where we have the
diagonal morphism (see Example~\ref{(co)diag-in-Loc})
\[
  \Delta\colon L \to L \oplus L
\]
as a closed localic map.

Id est, the image~$\Delta[L]$ is a closed sublocale in~$L \oplus L$.

\begin{lem}
  $\bigwedge \Delta[L] = \Delta(0_L)$.
\end{lem}

\begin{proof}
  Again, recall Example~\ref{(co)diag-in-Loc}\thinspace.

  $\subseteq$: The down-set $\Delta(0_L)$ is among elements of $\Delta[L]$; the
  infimum lies below all of~those.

  $\supseteq$: By the formula of the diagonal morphism,
  \[
    \Delta(0_L)
    = \left\{(x, y) \st x \wedge y \leq 0_L \right\}
    \subseteq \left\{(x, y) \st x \wedge y \leq a \right\}
    = \Delta(a)
  \]
  for every $a\in L$.
  From the property of the~infimum, $\Delta(0_L) \subseteq \bigwedge
  \Delta[L]$.
\end{proof}

Combining the above lemma with the~formula of~the sublocalic closure
\[
  \overline{S} = \left\uparrow \bigwedge \right. S,
\]
Isbell's condition reduces to
\[
  \Delta[L] = \left\uparrow \Delta(0_L),\right.
\]
which---after denoting the crucial element $\Delta(0_L)$ by $d_L$---produces
the final

\begin{framed}
  \begin{df}[I-Haus]
    A locale $L$ is considered \emph{strongly Hausdorff\/} (equally
    \emph{I-Hausdorff}: short for \emph{Isbell-Hausdorff}) iff
    \[
      \Delta[L] = \uparr d_L
    \]
  \end{df}
\end{framed}

\section{The weak from the strong one}

The current chapter is a~treatise of the relationship between strongly and
weakly Hausdorff property.
As their names suggest: we will make effort%
\footnote{hopefully joyful}
to prove the~following

\begin{thm}
  The (I-Haus) implies the (DS-Haus).
\end{thm}
(The strong version implies the weak one; hence the adjectives.)

However, at first, it may seem to be of some advantage to introduce several
technical lemmas.

\begin{lem}
  Suppose a join-basis $\mathcal{B}$ of a locale $L$ and a relation $R
  \subseteq L \times L$ possessing property
  \[
    \forall a, b\in L \; \forall c \in \mathcal{B}: \quad aRb \Rightarrow (a
    \wedge c)R(b \wedge c).
  \]
  Such a premise cause an element $s\in L$ to be $R$-saturated if and only if
  \[
    aRb \qquad \Rightarrow \qquad (a \leq s \; \Leftrightarrow \; b \leq s).
  \]
\end{lem}
\begin{proof}
  $\Rightarrow$: Trivial (put $c := 1$ in the definition of saturatedness).

  $\Leftarrow$: For a general $c\in L$ there are $c_i\in \mathcal{B}$ holding
  $c = \bigvee_{i\in J} c_i$.
  Following from frame distributivity, the~definition of~$\bigvee$ and the
  special property of~$R$, we get
  \begin{align*}
    a \wedge c \leq s \quad
    &\equiv a \wedge \bigvee_{i\in J} c_i \leq s \\
    &\equiv \bigvee_{i\in J} (a \wedge c_i) \leq s \\
    &\equiv \forall i\in J: a \wedge c_i \leq s \\
    &\equiv \forall i\in J: b \wedge c_i \leq s \\
    &\equiv \bigvee_{i\in J} (b \wedge c_i) \leq s \\
    &\equiv b \wedge \bigvee_{i\in J} c_i \leq s
    \equiv \quad b \wedge c \leq s;
  \end{align*}
  quite easily done.
\end{proof}

\begin{rem} \label{satur-def-eq}
  If additionally $aRb \Rightarrow a \leq b$, the previous alternative can be
  rewritten to
  \[
    aRb \quad \Rightarrow \quad (a \leq s \quad \Rightarrow \quad b \leq s)
  \]
  or even simpler
  \[
    aRb \; \& \; a \leq s \quad \Longrightarrow \quad b \leq s.
  \]

  Recall the products in {\bf Loc} (viz Preliminaries).
  The relation $R$ used in its definition satisfy both aforementioned
  conditions;
  whence, in case of binary products, $R$-saturated sets are precisely the
  down-sets $U \subseteq L_1 \times L_2$ fulfilling 
  \[
    \textstyle
    \{ (x_i, y) \st i \in J \} \subseteq U \quad \Longrightarrow \quad
    \left(\bigvee_{i\in J} x_i, y\right)\in U
  \]
  and
  \[
    \textstyle
    \{ (x, y_i) \st i \in J \} \subseteq U \quad \Longrightarrow \quad \left(x,
    \bigvee_{i\in J} y_i\right)\in U.
  \]
  Again, for further specification have a look at Picado and Pultr
  \cite{picado-pultr12}.
\end{rem}

\begin{lem} \label{downsets-satur}
  The down-set
  \[
    U = \left\downarrow(a, a \wedge b)\right. \cup \left\downarrow(a \wedge b, b)\right. \cup \n
  \]
  is saturated in $L \times L$ for arbitrary $a, b\in L$.
\end{lem}
\begin{proof}
  We are going to verify the saturation from~\ref{satur-def-eq}\thinspace.
  First of all, let us have $(x_i, y)\in U$ for $i\in J$.

  \underline{Case $y = 0$}:
  obviously, $(\bigvee_{i\in J} x_i, y) = (\bigvee_{i\in J} x_i, 0)\in \n
  \subseteq U$.

  \underline{Case $y \ne 0$ and $y \leq a \wedge b$}:
  then $x_i \leq a$ anyway, that is, $(\bigvee_{i\in J} x_i, y) \in
  \left\downarrow(a, a \wedge b)\right.$, which is a subset of $U$.

  \underline{Case $y \ne 0$ and $y \not\leq a \wedge b$}:
  in this case, inevitably, $y \leq b$ as well as $x_i \leq a \wedge b$; and
  analogously, again $(\bigvee_{i\in J} x_i, y) \in \left\downarrow(a \wedge b,
  b)\right. \subseteq U$.

  The $(x, \bigvee_{i\in J} y_i)$ by symmetry.
\end{proof}

\begin{lem} \label{meets-in-satur}
  In case of a strongly Hausdorff locale $L$, we have
  \[
    \Delta(\nabla(U)) = U
  \]
  for any saturated $U \supseteq d_L$.
\end{lem}
\begin{proof}
  By (I-Haus), every saturated $U\in \left\uparrow d_L \right. = \Delta[L]$ is an image
  $\Delta(a)$ for some $a\in L$.
  Let us write $\delta(U)$ for the mentioned $a$.
  In other words, $\Delta\delta(U) = U$.

  Remind yourselves of~Example~\ref{(co)diag-in-Loc}\thinspace.
  Notice $\Delta^*(a \oplus a) = a$;
  thence, the codiagonal~$\nabla$ is onto, and consequently, an epimorphism.
  Furthemore, using a standard property of~Galois adjoints $\nabla \Delta
  \nabla = \nabla$,
  we get
  \[
    \nabla \Delta = id.
  \]

  Joining the two observations results in
  \[
    \nabla (U) = \nabla (\Delta\delta (U)) = (\nabla \Delta)\delta (U) = \delta(U),
  \]
  which produce the desired $U = \Delta \delta (U) = \Delta \nabla (U)$.
\end{proof}

\begin{lem} \label{meets-in-satur}
  Let $L$ be an I-Hausdorff locale.
  Then, one obtains the~implication
  \[
    \left( a \wedge b, a \wedge b \right) \in U
    \; \Rightarrow \;
    \left( a, b \right) \in U
  \]
  for all saturated $U \supseteq d_L$.
\end{lem}
\begin{proof}
  Whenever $(a \wedge b, a \wedge b)\in U$, we see from the formula for
  $\nabla$ that
  \[
    a \wedge b \leq \bigvee \left\{ x \st (x, x) \in U \right\} = \nabla(U).
  \]
  Having in mind the formula for $\Delta$, immediately $(a, b) \in \Delta( \nabla(U) ) = U$
  (the final equality by the previous lemma).
\end{proof}

\begin{rem}
  The last and the penultimate implication actually hold in forms
  of~equivalences.
  Nevertheless, the reverse implications will be spared since we will make no
  use of them.
  A curious reader may survey Section~2 of Chapter~V in~Picado
  and~Pultr~\cite{picado-pultr12}.
\end{rem}

At last, it is high time to prove the promised theorem.
All the more, a~slightly stronger version of the (DS-Haus) definition
\[
  a \vee b \not\in \left\{a, b\right\} \qquad \Rightarrow \qquad \exists
  u\not\leq b, v\not\leq a: \quad u \wedge v = 0, \quad \boxed{\left(u,
  v\right) \leq \left(a, b\right)}
\]
will even be acquired.

In order to refresh:

\begin{thm} \label{IHaus->DSHaus}
  The (I-Haus) implies the (DS-Haus).
\end{thm}

\begin{proof}
  For a contradiction: let there be an I-Hausdorff $L$, which is not weakly
  Hausdorff.
  That means, we have $a, b$ with $a \not\leq b, \, b \not\leq a$ and such that
  \[
    u \wedge v = 0 \quad \& \quad \left(u, v\right) \leq \left(a, b\right)
    \qquad \Longrightarrow \qquad
    u \leq b \; \textrm{ or } \; v \leq a.
  \]
  Especially, for the down-set $U$ taken from~\ref{downsets-satur} we have
  \[
    d_L \cap (a \oplus b) \subseteq U.
  \]

  Following from~\ref{meets-in-satur}\thinspace, the saturated $\left(a \wedge
  b\right) \oplus \left(a \wedge b\right) \vee d_L \supseteq d_L$ contains $(a,
  b)$ (since it trivially contains $\left( a \wedge b, a \wedge b \right)$).
  Hence,
  \begin{align*}
    (a, b) &\in (a \oplus b) \; \cap \; ((a \wedge b) \oplus (a \wedge b) \vee d_L) \\
           &\subseteq ((a \wedge b) \oplus (a \wedge b)) \; \vee \; ((a \oplus b) \cap d_L) \\
           &\subseteq ((a \wedge b) \oplus (a \wedge b)) \; \vee \; U \\
           &= U.
  \end{align*}
  The first inclusion using~distributivity and the fact that $(a \wedge b)
  \oplus (a \wedge b) \subseteq a \oplus b$;
  the latter one by the observation from the beginning of this proof.
  The equality is seen from the~definition of~$U$:
  to be more specific, from $(a \wedge b) \oplus (a \wedge b) \subseteq U$.

  This leads to the long sought contradiction:
  $(a, b)\not\in \n$ since neither $a$ nor $b$ may equal $0$ (by $a \not\leq b,
  \, b \not\leq a$).
  Thus, the only options available are either $a \leq a \wedge b$ or $b \leq a
  \wedge b$.
  However, those enforce $a \leq b$ and $b \leq a$ respectively,
  a~contradiction. 
\end{proof}

\chapter{Weakly and strongly Hausdorff}

In this chapter we will consider point-free analogies of the $T_2$
axiom---a~separation not quite easy to imitate in frames.

\section{The $T_2$ axiom}

\begin{framed}
  \begin{df}[$T_2$]
    A topological space $(X, \tau)$ is called \emph{Hausdorff\/} (or
    a~\emph{$T_2$-space\/}) if
    \begin{center} \it
      for any $x \ne y$ from~$X$ there are disjoint $U$ and $V$ with $U\owns x$
      and $V\owns y$.
    \end{center}
  \end{df}
\end{framed}

\begin{rem} \label{T2->T1}
  Trivially, every $T_2$-space is $T_1$.
\end{rem}

\section{The Dowker-Strauss approach}

In 1972 Dowker and Strauss \cite{ds72} introduced the following condition:

\begin{framed}
  \begin{df}
    A locale $L$ satisfies the $S_2'$ axiom if
    \[
      (S_2') \qquad
      a, b \ne 1 \text{ and } a \vee b = 1 \quad \Rightarrow \quad \exists
      u\not\leq a, v\not\leq b, \quad u \wedge v = 0.
    \]
  \end{df}
\end{framed}

\begin{prop} \label{Haus->S2'}
  For every Hausdorff space $X$ the corresponding locale $\Omega(X)$ satisfies
  $S_2'$.
\end{prop}

\begin{proof}
  Open sets $A, B \ne X$ with $A \cup B = X$ contain $x\in X\setminus A =
  B\setminus A$ and $y\in X\setminus B = A\setminus B$.
  Particularly, we observe $x \ne y$ and, on top of that, receive the disjoint
  $U$ and $V$ from the $T_2$-definition.
  Evidently, $U\not\subseteq A$ (because of $x$) and $V\not\subseteq B$
  (because of $y$).
\end{proof}

\begin{rem}
  The reverse implication requires, in addition, the subfitness.
  For the~proof see Proposition 7 of~Dowker and Strauss
  \cite{ds72}.
\end{rem}

The $S_2'$ axiom is often replaced by a~somewhat stronger condition:

\begin{framed}
  \begin{df}[DS-Haus]
    A locale $L$ is considered to be \emph{weakly Hausdorff\/} (also
    \emph{DS-Hausdorff} as in \emph{Dowker-Straus-Hausdorff}) if
    \[
      a \vee b \not\in \left\{a, b\right\} \qquad \Rightarrow \qquad \exists
      u\not\leq a, v\not\leq b, \quad u \wedge v = 0.
    \]
  \end{df}
\end{framed}

\section{Isbell's approach}

Here is a~characterization of~the classical~$T_2$:

\begin{prop}
  A topological space $X$ is Hausdorff space if and only if the~diagonal
  $\Delta = \left\{(x, x) \st x\in X \right\}$ is closed in the product
  $X\times X$.
\end{prop}

\begin{proof}
  $\Rightarrow:$ Any element $(x, y)\not\in \Delta$, that means $x \ne y$, is
  separable from $\Delta$:
  namely, by~the~open $p_1^{-1}[U] \cap p_2^{-1}[V]$ with non-intersecting $U$
  and $V$ from the definition of~$(T_2)$.
  Hence, $(x, y)\not\in \overline{\Delta}$ concludes to $\Delta =
  \overline{\Delta}$, which is closed.

  $\Leftarrow:$ Choose $x \ne y$, in other words, $(x, y)\not\in \Delta$.
  Then $(x, y)$ has an open basic neighbourhood $U_1\times U_2$ disjoint from
  $\Delta$.
  That is, $U_1 \cap U_2 = \none$.
\end{proof}

This was used by Isbell \cite{isbell72} for the point-free counterpart of the
$T_2$ axiom:
Hausdorff locales are defined as those in which the diagonal (see
page~\pageref{codiag-in-Frm}\thinspace)
\[
  \Delta\colon L \to L \oplus L
\]
is a closed localic map.

\begin{lem}
  $\Delta$ is closed iff $\Delta[L] = \left\uparrow \Delta(0)\right.$.
\end{lem}
\begin{proof}
  $\Rightarrow$:
  By~\ref{prop:closed-map-char} on
  page~\pageref{prop:closed-map-char}\thinspace, $\Delta[L] =
  \Delta[\left\uparrow 0 \right.] = \left\uparrow \Delta(0)\right.$.

  $\Leftarrow$:
  We will prove the condition from~\ref{prop:closed-map-char}\thinspace, namely
  \[
    \Delta[ \left\uparrow a\right. ] = \left\uparrow \Delta(a).\right.
  \]

  The inclusion $\subseteq$ holds by~\ref{cor:delta-of-upsets} (viz
  page~\pageref{cor:delta-of-upsets}\thinspace).

  On the other hand, let $U\in \left\uparrow \Delta(a).\right.$
  Using the monotonicity of~$\Delta$, we see that $\left\uparrow
  \Delta(a)\right. \subseteq \left\uparrow \Delta(0).\right.$
  Therefore, we have $U\in \left\uparrow \Delta(0)\right. = \Delta[L]$ and
  hence $U = \Delta(b)$ for~some $b\in L$.
  Using Lemma~\ref{lem:delta-monotone} on page~\pageref{lem:delta-monotone} we
  get $a \le b$, and finally, $U = \Delta(b) \in \Delta[ \left\uparrow a\right.
  ]$.
\end{proof}

Thus---having denoted the crucial element $\Delta(0)$ by $d_L$---the condition
proposed by Isbell reduces~to

\begin{framed}
  \begin{df}[I-Haus]
    A locale $L$ is \emph{strongly Hausdorff\/} (or \emph{I-Hausdorff}: short
    for \emph{Isbell-Hausdorff}) iff
    \[
      \Delta[L] = \dL.
    \]
  \end{df}
\end{framed}

\begin{rem} \label{rem:delta-restriction}
  By~\ref{cor:delta-of-upsets} it always holds that $\Delta[L] \subseteq \dL$
  and hence we are allowed to (correctly) define
  \[
    \overline{\Delta}
  \]
  as the codomain restriction of~$\Delta$ to~$L \to \dL$.
\end{rem}

\begin{thm} \label{IHaus->DSHaus}
  (I-Haus) implies (DS-Haus).
\end{thm}

\begin{lem} \label{lem:delta-nabla=id}
  In case of a strongly Hausdorff locale $L$, we have
  \[
    \Delta(\nabla(U)) = U
  \]
  for any saturated $U \supseteq d_L$.
\end{lem}
\begin{proof}
  By (I-Haus), every saturated $U\in \dL = \Delta[L]$ is an image $\Delta(a)$
  for some $a\in L$.
  Let us write $\delta(U)$ for this $a$.
  In other words, $\Delta\delta(U) = U$.

  Recall the frame codiagonal on~page~\pageref{codiag-in-Frm}\thinspace.
  It is an epimorphism and hence---as~we have $\nabla \Delta \nabla = \nabla$
  by Fact~\ref{lrl=l} on page~\pageref{lrl=l}\thinspace---we also get $\nabla
  \Delta = id$.

  Joining these two observations results in
  \[
    \nabla (U) = \nabla (\Delta\delta (U)) = (\nabla \Delta)\delta (U) =
    \delta(U),
  \]
  which produce the desired $U = \Delta \delta (U) = \Delta \nabla (U)$.
\end{proof}

\begin{lem} \label{lem:mapping->IHaus}
  Let there be a~mapping $\alpha\colon L \to \dL$ such that
  \[
    \alpha \nabla = (U \mapsto U \vee d_L).
  \]
  Then $L$ is strongly Hausdorff.
\end{lem}
\begin{proof}
  Let $\gamma\colon L \oplus L \to \dL$ stand for $(U \mapsto U \vee d_L)$.
  Because $\gamma$ is onto, so is~$\alpha$.

  Furthermore, $\nabla$~and~$\gamma$ are frame homomorphism (using asociativity
  for joins and~distributivity for finite meets);
  and since $\nabla$ is onto, $\alpha$ is also a~frame homomorphism (recall
  \ref{lem:onto->homo} on page~\pageref{lem:onto->homo}\thinspace).
  Consequently, $\alpha$ is an~epimorphism.

  Since $\nabla \Delta (0) \le 0$ (from Fact~\ref{lr-id-rl} on
  page~\pageref{lr-id-rl}\thinspace), we have
  \[
    \nabla \alpha \nabla(U)
    = \nabla \gamma(U)
    = \nabla (U \vee d_L)
    = \nabla (U) \vee \nabla \Delta(0)
    = \nabla (U),
  \]
  that is, $\nabla \alpha = id$.
  By~\ref{lem:epi->iso} on page~\pageref{lem:epi->iso}\thinspace, $\alpha$ is
  an isomorphism and it is inverse to~the~$\nabla$ restricted~to $\dL \to L$.
  Consequently, by the uniqueness of~adjoints, $\alpha$ has to be
  the~restriction~$\overline{\Delta}$
  from~\ref{rem:delta-restriction}\thinspace.
  Finally, $\Delta[L] = \overline{\Delta}[L] = \alpha[L] = \dL$.
\end{proof}

\begin{prop} \label{meets-in-satur}
  A~locale~$L$ is I-Hausdorff iff one has the~implication
  \[
    \left( a \wedge b, a \wedge b \right) \in U
    \; \Rightarrow \;
    \left( a, b \right) \in U
  \]
  for all saturated $U \supseteq d_L$.
\end{prop}
\begin{proof}
  $\Rightarrow$:
  Whenever $(a \wedge b, a \wedge b)\in U$, we see from the formula for
  $\nabla$ that
  \[
    a \wedge b \leq \bigvee \left\{ x \st (x, x) \in U \right\} = \nabla(U).
  \]
  Having in mind the formula for $\Delta$, immediately $(a, b) \in \Delta(
  \nabla(U) ) = U$
  (the final equality by Lemma~\ref{lem:delta-nabla=id}\thinspace).

  $\Leftarrow$:
  Let the condition hold and~let $a_i\in L$ for $i\in J$. 
  Since for any~saturated down-set~$U \supseteq d_L$ we have $(a_i, a_i) \in U
  \Rightarrow (a_i \wedge a_j, a_i \wedge a_j) \in U$, we also obtain
  the~implication
  \begin{equation} \label{eq:IHaus-impl}
    (a_i, a_i) \in U \Longrightarrow \{ (a_i, a_j) \st j\in J \} \subseteq U
    %\tag{\arabic{chapter}.\thesection}
  \end{equation}

  We will prove the condition from~\ref{lem:mapping->IHaus}\, for
  $\alpha := (a \mapsto (a \oplus a) \vee d_L)$.
  Using
  \[
    \bigvee_{i\in J} a_i \oplus \bigvee_{j\in J} a_j
    = \iota_1 \left( \bigvee_{i\in J} a_i \right) \wedge \iota_2 \left(
    \bigvee_{j\in J} a_j \right)
    = \bigvee_{i, j\in J} \left( \iota_1(a_i) \wedge \iota_2(a_j) \right)
    = \bigvee_{i, j\in J} \left( a_i \oplus a_j \right)
  \]
  (the latter equality by double application of frame distributivity), we see
  that $\alpha$~preserves suprema:
  \begin{align*}
    \alpha \left( \bigvee_{i\in J} a_i \right)
    &= \left(\bigvee_{i, j\in J} \left( a_i \oplus a_j \right)\right) \vee d_L
    = \bigvee_{i\in J} \left(\left( a_i \oplus a_i \right) \vee d_L \right)
    = \bigvee_{i\in J} \alpha \left( a_i \right)
  \end{align*}
  (by~\eqref{eq:IHaus-impl} the second and the third expressions have identical
  upper bounds).
  Thus, both $\alpha$ and $\nabla$ preserve joins and~so does $\alpha \nabla$. 

  Furthermore, we have
  \[
    \alpha \nabla (a \oplus b)
    = \alpha (a \wedge b)
    = ((a \wedge b) \oplus (a \wedge b)) \vee d_L
    = (a \oplus b) \vee d_L,
  \]
  the last equality stipulated by the assumed implication.

  In~conclusion,
  \begin{flalign*}
    \alpha \nabla (U)
    &= \alpha \nabla \left(\bigvee \{a \oplus b \st (a, b)\in U\}\right)
    = \bigvee \{\alpha \nabla \left(a \oplus b \right) \st (a, b)\in U\} \\
    &= \bigvee \{a \oplus b \st (a, b)\in U\} \vee d_L
    = U \vee d_L,
  \end{flalign*}
  as the elements $a \oplus b$ generate $L \oplus L$ (see
  page~\pageref{a+b-gen}\thinspace).
\end{proof}

\begin{lem} \label{downsets-satur}
  The down-set
  \[
    U = \left\downarrow(a, a \wedge b)\right. \cup \left\downarrow(a \wedge b,
        b)\right. \cup \n
  \]
  is saturated in $L \times L$ for arbitrary $a, b\in L$.
\end{lem}
\begin{proof}
  Recall the saturatedness on~page~\pageref{df:satur}\thinspace.
  First of all, let us have $(x_i, y)\in U$ for $i\in J$.

  \underline{Case $y = 0$}:
  obviously, $(\bigvee_{i\in J} x_i, y) = (\bigvee_{i\in J} x_i, 0)\in \n$.

  \underline{Case $y \ne 0$ and $y \leq a \wedge b$}:
  then it must be that $x_i \leq a$ in~any case.
  Thus, $(\bigvee_{i\in J} x_i, y) \in \left\downarrow(a, a \wedge b).\right.$

  \underline{Case $y \ne 0$ and $y \not\leq a \wedge b$}:
  we have both $y \leq b$ and $x_i \leq a \wedge b$.
  Hence, again $(\bigvee_{i\in J} x_i, y) \in \left\downarrow(a \wedge b,
  b)\right. \subseteq U$.

  The $(x, \bigvee_{i\in J} y_i)$ by symmetry.
\end{proof}

Now we can prove the theorem.
In fact, we are going to imply a~stronger version of the (DS-Haus):
\[
  a \vee b \not\in \left\{a, b\right\} \qquad \Rightarrow \qquad \exists
  u\not\leq b, v\not\leq a: \quad u \wedge v = 0, \quad \boxed{\left(u,
  v\right) \leq \left(a, b\right)}
\]

\begin{proof}[Proof of~{\bf \ref{IHaus->DSHaus}}\thinspace]
  For a contradiction: let there be an I-Hausdorff $L$, which is not weakly
  Hausdorff.
  That means, we have $a, b$ with $a \not\leq b, \, b \not\leq a$ and such that
  \[
    u \wedge v = 0 \quad \& \quad \left(u, v\right) \leq \left(a, b\right)
    \qquad \Longrightarrow \qquad
    u \leq b \; \textrm{ or } \; v \leq a.
  \]
  Especially, for the down-set $U$ taken from~\ref{downsets-satur} we have
  \begin{equation} \label{eq:dl-aoplusb}
    d_L \cap (a \oplus b) \subseteq U.
  \end{equation}

  Following from~\ref{meets-in-satur}\thinspace, the saturated $(\left(a \wedge
  b\right) \oplus \left(a \wedge b\right)) \vee d_L \supseteq d_L$
  contains $(a, b)$ as it trivially contains $\left( a \wedge b, a \wedge b
  \right)$.
  Hence,
  \begin{align*}
    (a, b) &\in (a \oplus b) \; \cap \; (((a \wedge b) \oplus (a \wedge b))
            \vee d_L) \\
           &= ((a \wedge b) \oplus (a \wedge b)) \; \vee \; ((a \oplus b) \cap
               d_L) \\
           &\subseteq ((a \wedge b) \oplus (a \wedge b)) \; \vee \; U \\
           &= U.
  \end{align*}
  The first equality using~distributivity and the fact that $(a \wedge b)
  \oplus (a \wedge b) \subseteq a \oplus b$;
  the inclusion by~\ref{eq:dl-aoplusb}\thinspace;
  the final equality from $(a \wedge b) \oplus (a \wedge b) \subseteq U$.

  This leads to the contradiction:
  $(a, b)\not\in \n$ since neither $a$ nor $b$ may equal $0$ (by $a \not\leq b,
  \, b \not\leq a$).
  Thus, the only options available are either $a \leq a \wedge b$ or $b \leq a
  \wedge b$.
  However, those stipulate $a \leq b$ or $b \leq a$ respectively,
  a~contradiction. 
\end{proof}

\chapter{Regularity}

\section{The $T_3$ axiom}

\begin{framed}
  \begin{df}[$T_3$]
    We will consider a topological space $X$ to be \emph{regular\/} (or
    satisfying the~\emph{$T_3$-axiom\/}) when
    \begin{center} \it
      for any $x\in X$ and any closed $A \not\owns x$ there are open $V_1\owns
      x$ and $V_2\supseteq A$ disjoint from~one another.
    \end{center}
  \end{df}
\end{framed}

Even though the definition concerns points, it may be reformulated without
them.

\begin{prop} \label{reg-char}
  A space $X$ is regular iff each $U\in \Omega(X)$ can be depicted as
  \[
    U = \bigcup \{V\in \Omega(X) \st \overline{V} \subseteq U\}.
  \]
\end{prop}
\begin{proof}
  The inclusion~$\supseteq$ is obvious.

  $\Rightarrow$:
  Choose a $x\in U$.
  Since $A := X\setminus U \not\owns x$ is closed, the existence
  of~$V_1$~and~$V_2$ from $(T_3)$ is stipulated.
  The disjointedness $V_1 \cap V_2 = \none$ leads to $V_1\subseteq X\setminus
  V_2$, which is a closed set.
  Therefore, $\overline{V_1}\subseteq X\setminus V_2$.

  Furthermore, the relation $A\subseteq V_2$ is equivalent to the relation
  $X\setminus V_2\subseteq X \setminus A = U$.
  Hence, for $V_x := V_1$ one has $\overline{V_x}\subseteq X\setminus V_2\subseteq
  U$.
  Such a~system $\{V_x \st x\in U\}$ constitutes a~subset cover of $U$, and
  consequently, finishes the proof of~the~inclusion $\subseteq$.

  $\Leftarrow$:
  With $U := X\setminus A$ take $V_x$ from above and set $V_1 := V_x$ and $V_2
  := X\setminus \overline{V_x}$.
\end{proof}

Now we will eliminate the~closure in the~formula:

\begin{lem}
  $\overline{V}\subseteq U \quad \equiv \quad \exists W\in \Omega(X)\colon W \cap V =
  \none \;\; \& \;\; W \cup U = X$.
\end{lem}
\begin{proof}
  $\Rightarrow$:
  Take $W := X\setminus \overline{V}$.

  $\Leftarrow$:
  Recall closure's definition.
  Suppose $z\in \overline{V}\setminus U$.
  Since $W \cup U = X$, the $z$ must lie in the $W$;
  howbeit, the $W$ does not intersect the $V$---forcing $z\not\in
  \overline{V}$.
\end{proof}

Without loss of generality, the $W$ might be replaced by~the
pseudocomplement~$V^*$.
Explicitly, the $X\setminus \overline{V}$ in $\Omega(X)$.%

\begin{note}
Pseudocomplements always exist in frames:
setting
\[
  a^* := \bigvee \{x \st x \wedge a = 0\},
\]
we have
\[
  a^* \wedge a = \bigvee \{x \st x \wedge a = 0\} \wedge a = \bigvee \{x \wedge
  a \st x \wedge a = 0\} = 0
\]
(the second equality from frame distributivity).
Evidently, the implication
\[
  x \wedge a = 0 \Rightarrow x \le a^*
\]
also holds.
\end{note}

\section{Regular locales}

\begin{framed}
  \begin{nota}[$\rb$]
    \[
      V \rb U \quad \equiv \quad V^* \vee U = 1,
    \]
    which is referred to by pronouncing \emph{``V is rather below U''\/}.
  \end{nota}
\end{framed}

The next lemma will be advantageous property for the rest of this chapter.
\begin{lem} \label{rb->leq}
  $a \rb b \Rightarrow a \leq b$.
\end{lem}
\begin{proof}
  By~distributivity,
  \[
    a = a \wedge 1 = a \wedge (a^* \vee b) = (a \wedge a^*) \vee (a \wedge b) =
    0 \vee (a \wedge b) = a \wedge b. \qedhere
  \]
\end{proof}

Possessing $\rb$ notation, we are able to adopt the characterization
\ref{reg-char} by~defining regularity as~following
\begin{framed}
  \begin{df}[Reg]
    A locale is called \emph{regular\/} if and only if
    \[
      a = \bigvee \{x \st x \rb a\}
    \]
    for all its elements $a$.
  \end{df}
\end{framed}

Thus we have that
\begin{center} 
  \emph{a topological space $X$ is regular iff its $\Omega(X)$ is regular.\/}
\end{center}

\section{Strength of regularity}

\begin{lem} \label{oplus-vee-distrib}
  ~
  \begin{enumerate}
  \item $\bigvee_{i\in J} \left(a_i \oplus b\right) = \left(\bigvee_{i\in J}
    a_i \right) \oplus b$.
  \item $\bigvee_{i\in J} \left(a \oplus b_i\right) = a \oplus
    \left(\bigvee_{i\in J} b_i \right)$.
  \item $\left(\bigvee_{i\in J} a_i\right) \oplus \left(\bigvee_{j\in J}
    b_j\right) = \bigvee_{i, j\in J} \left(a_i \oplus b_j\right)$.
  \end{enumerate}
\end{lem}
\begin{proof}
  (i):
  Recall the corollary on page~\pageref{oplus-iota}\thinspace.
  We have
  \begin{align*}
    \left(\bigvee_{i\in J} a_i \right) \oplus b
    &= \iota_1\left(\bigvee_{i\in J} a_i \right) \wedge \iota_2(b)
    = \left(\bigvee_{i\in J} \iota_1(a_i) \right) \wedge \iota_2(b) \\
    &= \bigvee_{i\in J} \left(\iota_1(a_i) \wedge \iota_2(b) \right)
    = \bigvee_{i\in J} \left(a_i \oplus b\right)
  \end{align*}
  as frame homomorphisms $\iota_i$ preserve suprema (for the second equality)
  and as the binary meet commutes with joins in frames (for the third
  equality).

  (ii):
  Analogously by symmetry.

  (iii):
  By consecutive application of (i) and (ii).
\end{proof}

\begin{lem}
  For a general locale $L$ and any of its saturated $U\in \left\uparrow
  d_L\right.$
  \[
    (a \wedge b, a \wedge b) \in U \qquad \Longrightarrow \qquad \forall x \rb
    a, \, y \rb b: \; (x, y) \in U
  \]
\end{lem}
\begin{proof}
  Beginning with evident $(x, y)\in x \oplus y$, we modify
  \begin{align*}
    x \oplus y &= (x \wedge (y^* \vee b)) \oplus (y \wedge (x^* \vee a)) \\
               &= (\; (x \wedge y^*) \; \vee \; (x \wedge b) \; ) \oplus (\; (y
    \wedge x^*) \; \vee \; (y \wedge a) \; )
  \end{align*}
  (by using ``rather-belowness'' and distributivity, respectively).

  Proceeding with (iii) of Lemma~\ref{oplus-vee-distrib}\thinspace,
  \begin{align*}
     \ldots = \; &(\; (x \wedge y^*) \oplus (y \wedge x^*) \; ) \vee 
            (\; (x \wedge y^*) \oplus (y \wedge a) \; ) \; \vee \\
            \vee \; &(\; (x \wedge b) \oplus (y \wedge x^*) \; ) \vee
            (\; (x \wedge b) \oplus (a \wedge y) \; ) \\
     \leq \; &(y^*\oplus y) \vee (y^*\oplus y) \vee (x\oplus x^*) \vee
                  (\; (a \wedge b)\oplus(a \wedge b) \;)
  \end{align*}
  where the last upper-bound is retrieved from~\ref{rb->leq}\thinspace.

  Besides, the ultimate expression is a subset of $U$:
  as $(x^*\oplus x), (y^*\oplus y)\subseteq d_L$ and $(a \wedge b)\oplus(a
  \wedge b)\subseteq U$ from the premise---all the way leading to $(x, y)\in
  U$.
\end{proof}

\begin{thm}
  (Reg) implies (I-Haus).
\end{thm}
\begin{proof}
  We will show that regular locales meet the condition
  from~\ref{meets-in-satur} on page~\pageref{meets-in-satur}\thinspace.
  Let $(a \wedge b, a \wedge b) \in U$.
  By the earlier lemma, we have $(x, y)\in U$ with $x \rb a$ and $y \rb b$.
  Recall the saturatedness on~page~\pageref{df:satur}\thinspace:
  as $U$ is saturated, we get
  \[
    \left\{ (x, y) \st x \rb a \right\} \subseteq U
    \Longrightarrow
    \left(\bigvee \{x \st x \rb a\}, y\right)\in U,
  \]
  or equally, using regularity,
  \[
    \left(a, y\right)\in U
  \]
  for all $y \rb b$.
  Wherefore, in the same manner: $(a, b) = \left(a, \bigvee \{y \st y \rb
  b\}\right)\in U$.
\end{proof}

By Theorem~\ref{IHaus->DSHaus} on page~\pageref{IHaus->DSHaus}\thinspace, we
get
\begin{cor}
  Regular locales are weakly Hausdorff.
\end{cor}

It will be useful to characterize (Reg) by a formula that resembles the subfit
definition:
\begin{lem} \label{reg-char}
  A locale is regular iff
  \[
    a \not\le b \qquad \Rightarrow \qquad \exists c: \quad a \vee c = 1 \quad
    \& \quad c^* \not\leq b
  \]
\end{lem}
\begin{proof}
  $\Rightarrow$:
  Suppose $a \not\le b$;
  that is, using the regularity, there is $x \rb a$ with $x \not\le b$ (otherwise
  $\bigvee \{x \st x \rb a\} \le b$).
  Let $c := x^*$.
  In other words, by $x \rb a$, one concludes $a \vee c = a \vee x^* = 1$.
  Additionally, $c^* \not\le b$: since otherwise from pseudocomplement property
  $x \le x^{**} = c^* \le b$, a contradiction.

  $\Leftarrow$:
  By~\ref{rb->leq} we always have the inequality $\bigvee \{x \st x \rb a\} \le
  a$.
  Hence, for contradiction assume $a \not\le \bigvee \{x \st x \rb a\}$.
  Then, there exists $c$ from the~premise.
  Specifically, $1 = a \vee c \le a \vee c^{**}$, which gives us $c^* \rb a$.
  Yet $c^* \not\le \bigvee \{x \st x \rb a\}$.
\end{proof}

\begin{thm}
  (Reg) implies (Sfit).
\end{thm}
\begin{proof}
  Given $a \not\le b$, we acquire---from Lemma~\ref{reg-char}\thinspace---an
  element $c$ such that $a \vee c = 1$ and $c^* \not\le b$.
  What is more, it satisfies $b \vee c \ne 1$.
  If not so then
  \[
    b = b \vee 0 = b \vee (c \wedge c^*) = (b \vee c) \wedge (b \vee c^*) = 1
    \wedge (b \vee c^*) = b \vee c^*,
  \]
  the third identity from distributivity; subsequently, this contradicts $c^*
  \le b$.
\end{proof}

\chapter{Complete regularity}

\section{The $T_{3\frac{1}{2}}$ axiom}

\begin{framed}
  \begin{df}[$T_{3\frac{1}{2}}$]
    A topological space $X$ is said to be \emph{completely regular\/} (or
    equally \emph{$T_{3\frac{1}{2}}$-space\/}) if
    \begin{center} \it
      for any $x\in X$ and any closed $A$ such that $x\not\in A$ there is a
      continuous function $\phi\colon X \to \I$ with $\phi(x) = 0$ and $\phi[A]
      = \{1\}$.
    \end{center}
  \end{df}
\end{framed}

Again, we introduce

\begin{framed}
  \begin{nota}[$\cb$]
    For elements $a, b$ of a locale $L$ let us write
    \[
      a \cb b
    \]
    and say \emph{``a is completely below b''\/} if for every $r\in \langle 0,
    1 \rangle \cap \Q$ there exist $a_r\in L$ satisfying
    \begin{align*}
      a_0 &= a, \\
      a_1 &= b, \\
      \qquad a_p &\rb a_q \quad \text{ for } p < q
    \end{align*}
  \end{nota}
\end{framed}

\begin{rem}
  Evidently, by~\ref{rb->leq} from page~\pageref{rb->leq}\thinspace, $a \cb b$
  implies $a \le b$.
\end{rem}

\begin{lem} \label{cb->continuous}
  Let $X$ be a~topological space and let $\{a_r\in \Omega(X) \st r\in \langle
  0, 1 \rangle \cap \Q\}$ be set witnessing $a \cb b$.
  Then function $f\colon X \to \I$ defined by
  \[
    f(x) := \inf \{r \st a_r \owns x\}
  \]
  is continuous.
\end{lem}
\begin{proof}
  TBD
\end{proof}

\begin{prop} \label{T3,5-char}
  A topological space $X$ is completely regular iff
  \[
    \forall U\in \Omega(X) \qquad U = \bigcup\{V\in \Omega(X) \st V \cb U\}
  \]
\end{prop}
\begin{proof}
  TBD
\end{proof}

In the analogous manner as in regular case:

\begin{framed}
  \begin{df}[CReg]
    A locale is called \emph{completely regular\/} if
    \[
      a = \bigvee \{x \st x \cb a\}
    \]
    for all its elements $a$.
  \end{df}
\end{framed}

\chapter{Normality}

\section{The $T_4$ axiom}

\begin{framed}
  \begin{df}[$T_4$]
    A topological space $X$ is \emph{normal\/} (or \emph{$T_4$-space\/}) if
    \begin{center} \it
      for every disjoint closed $A, B$ there exist disjoint open $U, V$ with
      $U\supseteq A, \, V\supseteq B$.
    \end{center}
  \end{df}
\end{framed}

\section{Normal locales}

The point-free counterpart to $T_4$ is straightforward:

\begin{framed}
  \begin{df}[Norm]
    A locale $L$ is \emph{normal\/} if
    \[
      a \vee b = 1 \qquad \Rightarrow \qquad \exists u, v, \quad u \wedge v =
      0 \text{ and } a \vee u = 1 = b \vee v.
    \]
  \end{df}
\end{framed}

\begin{rem}
  As $u \wedge v = 0 \; \equiv \; v \le u^*$, we may define (Norm) equally by
  \[
    a \vee b = 1 \qquad \Rightarrow \qquad \exists u, \quad a \vee u = 1 = b
    \vee u^*.
  \]
\end{rem}

\medskip

Because the definition of $(T_4)$ is virtually point-free, we obtain following

\begin{cor}
  A topology $X$ is normal iff the locale $\Omega(X)$ is normal.
\end{cor}

\begin{lem}
  The relation $\rb$ interpolates in~normal locales.
\end{lem}
\begin{proof}
  Suppose $a$ is rather below $b$.
  That is, $a^* \vee b = 1$, and using the normality,\thinspace%
  \footnote{more precisely, the subsequent remark}
  we get $u\in L$ such that
  $a^* \vee u = 1 = b \vee u^*$.
  In other words, $a \rb u \rb b$.
\end{proof}

Combined with Remark~\ref{cb-largest-interpol} on
page~\pageref{cb-largest-interpol}\thinspace, we form

\begin{cor}
  (Norm) implies $\rb \, = \, \cb$;
  hence, in case of normal locales, regularity coincides with complete
  regularity.
\end{cor}

\begin{thm}
  (Norm) $\&$ (Sfit) $\Rightarrow$ (Reg).
\end{thm}
\begin{proof}
  Take $a\in L$ and put $b := \bigvee \{x \st x \rb a\}$.
  For contradiction, suppose that $a\not\le b$.
  Using the subfitness, there is $c\in L$ fulfilling $a \vee c = 1 \ne b \vee
  c$.
  Applying the normality, we find $u\in L$ with $a \vee u^* = 1 = c \vee u$.
  Particularly, $u \rb a$, which leads to~$u \le b$.
  Finally, $b \vee c \ge u \vee c = 1$; contradicting the subfitness.
\end{proof}

Since $\rb$ equals $\cb$ in normal locales, we also deduce

\begin{cor}
  (Norm) $\&$ (Sfit) $\Rightarrow$ (CReg).
\end{cor}

\chapter{Normality}

\section{The $T_4$ axiom}

\begin{framed}
  \begin{df}[$T_4$]
    A topological space $X$ is \emph{normal\/} (or \emph{$T_4$-space\/}) if
    \begin{center} \it
      for every disjoint closed $A, B$ there exist disjoint open $U, V$ with
      $U\supseteq A, \, V\supseteq B$.
    \end{center}
  \end{df}
\end{framed}

\section{Normal locales}

The point-free counterpart to $T_4$ is straightforward:

\begin{framed}
  \begin{df}[Norm]
    A locale $L$ is \emph{normal\/} if
    \[
      a \vee b = 1 \qquad \Rightarrow \qquad \exists u, v, \quad u \wedge v =
      0 \text{ and } a \vee u = 1 = b \vee v.
    \]
  \end{df}
\end{framed}

\begin{rem}
  As $u \wedge v = 0 \; \equiv \; v \le u^*$, we may define (Norm) equally by
  \[
    a \vee b = 1 \qquad \Rightarrow \qquad \exists u, \quad a \vee u = 1 = b
    \vee u^*.
  \]
\end{rem}

\medskip

Because the definition of $(T_4)$ is virtually point-free, we obtain following

\begin{cor}
  A topology $X$ is normal iff the locale $\Omega(X)$ is normal.
\end{cor}

\begin{lem}
  The relation $\rb$ interpolates in~normal locales.
\end{lem}
\begin{proof}
  Suppose $a$ is rather below $b$.
  That is, $a^* \vee b = 1$, and using the normality,\thinspace%
  \footnote{more precisely, the subsequent remark}
  we get $u\in L$ such that
  $a^* \vee u = 1 = b \vee u^*$.
  In other words, $a \rb u \rb b$.
\end{proof}

Combined with Remark~\ref{cb-largest-interpol} on
page~\pageref{cb-largest-interpol}\thinspace, we form

\begin{cor}
  (Norm) implies $\rb \, = \, \cb$;
  hence, in case of normal locales, regularity coincides with complete
  regularity.
\end{cor}

\begin{thm}
  (Norm) $\&$ (Sfit) $\Rightarrow$ (Reg).
\end{thm}
\begin{proof}
  Take $a\in L$ and put $b := \bigvee \{x \st x \rb a\}$.
  For contradiction, suppose that $a\not\le b$.
  Using the subfitness, there is $c\in L$ fulfilling $a \vee c = 1 \ne b \vee
  c$.
  Applying the normality, we find $u\in L$ with $a \vee u^* = 1 = c \vee u$.
  Particularly, $u \rb a$, which leads to~$u \le b$.
  Finally, $b \vee c \ge u \vee c = 1$; contradicting the subfitness.
\end{proof}

Since $\rb$ equals $\cb$ in normal locales, we also deduce

\begin{cor}
  (Norm) $\&$ (Sfit) $\Rightarrow$ (CReg).
\end{cor}

\chapter{Compactification}

\section{The Čech-Stone compactification}

%\begin{framed}
%  \begin{df}[$\beta X$]
%    \emph{The Čech-Stone compactification\/} of a topological space $X$ is a
%    compact space
%    \[
%      \beta X
%    \]
%    with the property \ldots
%  \end{df}
%\end{framed}


% Ukázka použití některých konstrukcí LateXu (odkomentujte, chcete-li)
% %%% Ukázka použití některých konstrukcí LaTeXu

\subsection{Ukázka \LaTeX{}u}
\label{ssec:ukazka}

This short subsection serves as an~example of basic \LaTeX{} constructs,
which can be useful for writing a~thesis.

Let us start with lists:

\begin{itemize}
\item The logo of Matfyz is displayed in figure~\ref{fig:mff}.
\item This is subsection~\ref{ssec:ukazka}.
\item Citing literature~\cite{lamport94}.
\end{itemize}

Different kinds of dashes:
red-black (short),
pages 16--22 (middle),
$45-44$ (minus),
and this is --- as you could have expected --- a~sentence-level dash,
which is the longest.
(Note that we have follwed \verb|a| by a~tilde instead of a~space
to avoid line breaks at that place.)

\newtheorem{theorem}{Theorem}
\newtheorem*{define}{Definition}	% Definice nečíslujeme, proto "*"

\begin{define}
A~{\sl Tree} is a connected graph with no cycles.
\end{define}

\begin{theorem}
This theorem is false.
\end{theorem}

\begin{proof}
False theorems do not have proofs.
\end{proof}

\begin{figure}
	\centering
	\includegraphics[width=30mm]{../img/logo.eps}
	\caption{Logo of MFF UK}
	\label{fig:mff}
\end{figure}


\chapter*{Epilogue}
\addcontentsline{toc}{chapter}{Epilogue}


%%% Seznam použité literatury
%%% Seznam použité literatury je zpracován podle platných standardů. Povinnou citační
%%% normou pro bakalářskou práci je ISO 690. Jména časopisů lze uvádět zkráceně, ale jen
%%% v kodifikované podobě. Všechny použité zdroje a prameny musí být řádně citovány.

\def\bibname{Bibliography}
\begin{thebibliography}{99}
\addcontentsline{toc}{chapter}{\bibname}

%\bibitem{lamport94}
%  {\sc Lamport,} Leslie.
%  \emph{\LaTeX: A Document Preparation System}.
%  2. vydání.
%  Massachusetts: Addison Wesley, 1994.
%  ISBN 0-201-52983-1.

\bibitem{ds72}
  {\sc Dowker,} Clifford Hugh and
  {\sc Strauss,} Dona Papert.
  Separation axioms for frames.
  In: \emph{Topics in Topology}.
  1972,
  vol.~8,
  pp.~223-240.
  Proc. Colloq., Keszthely (Hungary), 1972.
  Colloq. Math. Soc. Janos Bolyai, Amsterdam: North-Holland, 1974.

\bibitem{isbell72}
  {\sc Isbell,} John Rolfe.
  Atomless parts of spaces.
  In: \emph{Mathematica Scandinavica}.
  1972,
  31:5-32.

\bibitem{johnstone-shu-hao}
  {\sc Johnstone,} Peter and
  {\sc Shu-Hao,} Sun.
  Weak products and Hausdorff locales.
  preprint.

\bibitem{paseka-smarda92}
  {\sc Paseka,} Jan and
  {\sc Šmarda,} Bohumil.
  $T_2$-frames and almost compact frames.
  In: \emph{Czechoslovak Math. Journal}.
  1992,
  vol.~42 (117),
  pp.~385-402.
  Praha: Czechoslovak Math. Journal, 1992.

\bibitem{picado-pultr12}
  {\sc Picado,} Jorge and
  {\sc Pultr,} Aleš.
  \emph{Frames and Locales: Topology without points}.
  2012 edition.
  Basel: Birkh\"{a}user, 2012.
  ISBN 978-3-0348-0153-9.

\bibitem{rosicky-smarda85}
  {\sc Paseka,} J. and
  {\sc Šmarda,} B.
  $T_1$-locales.
  In: \emph{Math. Proc. Cambr. Phil. Soc}.
  1985,
  vol.~98,
  pp.~81-86.

\end{thebibliography}


%%% Tabulky v bakalářské práci, existují-li.
%\chapwithtoc{List of Tables}

%%% Použité zkratky v bakalářské práci, existují-li, včetně jejich vysvětlení.
%\chapwithtoc{List of Abbreviations}

%%% Přílohy k bakalářské práci, existují-li (různé dodatky jako výpisy programů,
%%% diagramy apod.). Každá příloha musí být alespoň jednou odkazována z vlastního
%%% textu práce. Přílohy se číslují.
%\chapwithtoc{Attachments}

\openright
\end{document}
